%=============================================================================
% Chapter 2 Solutions
%=============================================================================
\newpage
\section*{Chapter 2 Solutions}
\addcontentsline{toc}{section}{Chapter 2 Solutions}

%-----------------------------------------------------------------------------
\subsection*{2.1 --- C to MIPS: Arithmetic}
\hypertarget{sol-2.1}{}
\markright{2.1}
\addcontentsline{toc}{subsection}{2.1 --- C to MIPS: Arithmetic}

\texttt{f = g + h - 5;} with \texttt{f}$=$\texttt{\$s0},
\texttt{g}$=$\texttt{\$s1}, \texttt{h}$=$\texttt{\$s2}.

\begin{verbatim}
    add  $s0, $s1, $s2     # f = g + h
    addi $s0, $s0, -5      # f = f - 5
\end{verbatim}
Two instructions (minimal).

%-----------------------------------------------------------------------------
\subsection*{2.2 --- MIPS to C: Two Adds}
\hypertarget{sol-2.2}{}
\markright{2.2}
\addcontentsline{toc}{subsection}{2.2 --- MIPS to C: Two Adds}

\begin{verbatim}
    add $t0, $s1, $s2      # t0 = g + h
    add $s0, $t0, $s3      # f = t0 + i
\end{verbatim}
C statement: \texttt{f = g + h + i;}

(Assuming \texttt{f}$=$\texttt{\$s0}, \texttt{g}$=$\texttt{\$s1},
\texttt{h}$=$\texttt{\$s2}, \texttt{i}$=$\texttt{\$s3}.)

%-----------------------------------------------------------------------------
\subsection*{2.3 --- C to MIPS: Array Access}
\hypertarget{sol-2.3}{}
\markright{2.3}
\addcontentsline{toc}{subsection}{2.3 --- C to MIPS: Array Access}

\texttt{B[8] = A[i] + A[j];} with 4-byte (word) elements.
\begin{verbatim}
    sll  $t0, $s3, 2       # $t0 = i * 4
    add  $t0, $s6, $t0     # $t0 = &A[i]
    lw   $t0, 0($t0)       # $t0 = A[i]
    sll  $t1, $s4, 2       # $t1 = j * 4
    add  $t1, $s6, $t1     # $t1 = &A[j]
    lw   $t1, 0($t1)       # $t1 = A[j]
    add  $t0, $t0, $t1     # $t0 = A[i] + A[j]
    sw   $t0, 32($s7)      # B[8] = $t0  (8*4=32)
\end{verbatim}

%-----------------------------------------------------------------------------
\subsection*{2.4 --- MIPS to C: Array Code}
\hypertarget{sol-2.4}{}
\markright{2.4}
\addcontentsline{toc}{subsection}{2.4 --- MIPS to C: Array Code}

Reading the assembly line by line:
\begin{enumerate}
  \item \texttt{sll \$t0, \$s0, 2} $\rightarrow$ \texttt{\$t0 = f * 4}
  \item \texttt{add \$t0, \$s6, \$t0} $\rightarrow$ \texttt{\$t0 = \&A[f]}
  \item \texttt{sll \$t1, \$s1, 2} $\rightarrow$ \texttt{\$t1 = g * 4}
  \item \texttt{add \$t1, \$s7, \$t1} $\rightarrow$ \texttt{\$t1 = \&B[g]}
  \item \texttt{lw \$s0, 0(\$t0)} $\rightarrow$ \texttt{f = A[f]}
  \item \texttt{addi \$t2, \$t0, 4} $\rightarrow$ \texttt{\$t2 = \&A[f] + 4
        = \&A[f+1]}
  \item \texttt{lw \$t0, 0(\$t2)} $\rightarrow$ \texttt{\$t0 = A[f+1]}
  \item \texttt{add \$t0, \$t0, \$s0} $\rightarrow$ \texttt{\$t0 = A[f+1] + f}
        (but f was overwritten in step 5 to A[f])
  \item \texttt{sw \$t0, 0(\$t1)} $\rightarrow$ \texttt{B[g] = \$t0}
\end{enumerate}

C code (using original f):
\begin{verbatim}
    B[g] = A[f] + A[f + 1];
\end{verbatim}
Note: the original value of \texttt{f} is destroyed (overwritten with
\texttt{A[f]}).

%-----------------------------------------------------------------------------
\subsection*{2.5 --- Byte Ordering: Endianness}
\hypertarget{sol-2.5}{}
\markright{2.5}
\addcontentsline{toc}{subsection}{2.5 --- Byte Ordering: Endianness}

Value \texttt{0xABCDEF12} stored at address 0:

\begin{center}
\begin{tabular}{lcccc}
\toprule
\textbf{Address} & 0 & 1 & 2 & 3 \\
\midrule
Big-endian    & \texttt{AB} & \texttt{CD} & \texttt{EF} & \texttt{12} \\
Little-endian & \texttt{12} & \texttt{EF} & \texttt{CD} & \texttt{AB} \\
\bottomrule
\end{tabular}
\end{center}

Big-endian stores the most significant byte at the lowest address.
Little-endian stores the least significant byte at the lowest address.

%-----------------------------------------------------------------------------
\subsection*{2.6 --- Hex to Decimal Conversion}
\hypertarget{sol-2.6}{}
\markright{2.6}
\addcontentsline{toc}{subsection}{2.6 --- Hex to Decimal Conversion}

\texttt{0xABCDEF12}:
\begin{align*}
  &= 10 \times 16^7 + 11 \times 16^6 + 12 \times 16^5 + 13 \times 16^4
     + 14 \times 16^3 + 15 \times 16^2 + 1 \times 16^1 + 2 \times 16^0 \\
  &= 2{,}882{,}400{,}018
\end{align*}

As a signed 32-bit integer (MSB is 1):
$2{,}882{,}400{,}018 - 2^{32} = -1{,}412{,}567{,}278$.

%-----------------------------------------------------------------------------
\subsection*{2.7 --- C to MIPS: 8-Byte Word Arrays}
\hypertarget{sol-2.7}{}
\markright{2.7}
\addcontentsline{toc}{subsection}{2.7 --- C to MIPS: 8-Byte Word Arrays}

Same as 2.3 but with 8-byte elements, so multiply index by 8 (shift left by
3):
\begin{verbatim}
    sll  $t0, $s3, 3       # $t0 = i * 8
    add  $t0, $s6, $t0     # $t0 = &A[i]
    lw   $t0, 0($t0)       # $t0 = A[i] (lower 32 bits)
    sll  $t1, $s4, 3       # $t1 = j * 8
    add  $t1, $s6, $t1     # $t1 = &A[j]
    lw   $t1, 0($t1)       # $t1 = A[j] (lower 32 bits)
    add  $t0, $t0, $t1     # $t0 = A[i] + A[j]
    sw   $t0, 64($s7)      # B[8] = $t0  (8*8=64)
\end{verbatim}
Note: for a full 64-bit value on a 32-bit MIPS, both the upper and lower
32 bits would need to be loaded and added. The solution above handles only
the lower word.

%-----------------------------------------------------------------------------
\subsection*{2.8 --- Instruction Format Fields}
\hypertarget{sol-2.8}{}
\markright{2.8}
\addcontentsline{toc}{subsection}{2.8 --- Instruction Format Fields}

\texttt{add \$s0, \$s0, \$s1} --- R-type:
\begin{center}
\begin{tabular}{cccccc}
\toprule
op & rs & rt & rd & shamt & funct \\
\midrule
0 & 16 (\$s0) & 17 (\$s1) & 16 (\$s0) & 0 & 32 (add) \\
\bottomrule
\end{tabular}
\end{center}

\texttt{sw \$t1, 0(\$s0)} --- I-type:
\begin{center}
\begin{tabular}{cccc}
\toprule
op & rs & rt & immediate \\
\midrule
43 (sw) & 16 (\$s0) & 9 (\$t1) & 0 \\
\bottomrule
\end{tabular}
\end{center}

%-----------------------------------------------------------------------------
\subsection*{2.9 --- MIPS Fields for Exercise 2.8}
\hypertarget{sol-2.9}{}
\markright{2.9}
\addcontentsline{toc}{subsection}{2.9 --- MIPS Fields for Exercise 2.8}

Binary encoding for each instruction:

\texttt{add \$s0, \$s0, \$s1}:
\begin{verbatim}
  000000 10000 10001 10000 00000 100000
  op=0   rs=16 rt=17  rd=16 sh=0  fn=32
\end{verbatim}
Hex: \texttt{0x02118020}

\texttt{sw \$t1, 0(\$s0)}:
\begin{verbatim}
  101011 10000 01001 0000000000000000
  op=43  rs=16 rt=9   imm=0
\end{verbatim}
Hex: \texttt{0xAE090000}

%-----------------------------------------------------------------------------
\subsection*{2.10 --- Register Arithmetic and Overflow}
\hypertarget{sol-2.10}{}
\markright{2.10}
\addcontentsline{toc}{subsection}{2.10 --- Register Arithmetic and Overflow}

\texttt{\$s0 = 0x80000000} $= -2{,}147{,}483{,}648$, \quad
\texttt{\$s1 = 0xD0000000} $= -805{,}306{,}368$.

\phantomsection\addcontentsline{toc}{subsubsection}{2.10.1}
\medskip\noindent\textbf{2.10.1} \texttt{add \$t0, \$s0, \$s1}:
$\texttt{0x80000000} + \texttt{0xD0000000} = \texttt{0x(1)50000000}$.
Lower 32 bits: \texttt{\$t0 = 0x50000000}.

\noindent\textbf{2.10.2}
Both operands are negative; the true sum is
$-2{,}147{,}483{,}648 + (-805{,}306{,}368) = -2{,}952{,}789{,}016$,
which is below $-2^{31}$.  The result \texttt{0x50000000} is positive.
\textbf{Overflow has occurred.}

\phantomsection\addcontentsline{toc}{subsubsection}{2.10.3}
\medskip\noindent\textbf{2.10.3} \texttt{sub \$t0, \$s0, \$s1}:
$\texttt{0x80000000} - \texttt{0xD0000000} = \texttt{0x80000000}
+ \texttt{0x30000000} = \texttt{0xB0000000}$.
\texttt{\$t0 = 0xB0000000} $= -1{,}342{,}177{,}280$.

\noindent\textbf{2.10.4}
True result: $-2{,}147{,}483{,}648 - (-805{,}306{,}368)
= -1{,}342{,}177{,}280$, which fits in 32 bits.
\textbf{No overflow.}

\phantomsection\addcontentsline{toc}{subsubsection}{2.10.5}
\medskip\noindent\textbf{2.10.5}
First: \texttt{add \$t0, \$s0, \$s1}
$\rightarrow$ \texttt{\$t0 = 0x50000000} (overflow already).
Then: \texttt{add \$t0, \$t0, \$s0}
$= \texttt{0x50000000} + \texttt{0x80000000} = \texttt{0xD0000000}$.

\noindent\textbf{2.10.6}
The first add already overflowed. The final value \texttt{0xD0000000}
happens to equal \texttt{\$s1}, but the computation is incorrect.
\textbf{Overflow occurred} in the first instruction.

%-----------------------------------------------------------------------------
\subsection*{2.11 --- Overflow Ranges}
\hypertarget{sol-2.11}{}
\markright{2.11}
\addcontentsline{toc}{subsection}{2.11 --- Overflow Ranges}

\texttt{\$s1 = 128}.

\noindent\textbf{2.11.1} \texttt{add \$s0, \$s0, \$s1}: overflow when
$\texttt{\$s0} + 128 > 2^{31}-1$, i.e.\ $\texttt{\$s0} > 2^{31}-129
= 2{,}147{,}483{,}519$.
Range: $\texttt{\$s0} \in [2{,}147{,}483{,}520,\; 2{,}147{,}483{,}647]$.

\noindent\textbf{2.11.2} \texttt{sub \$s0, \$s0, \$s1}: overflow when
$\texttt{\$s0} - 128 < -2^{31}$, i.e.\ $\texttt{\$s0} < -2^{31}+128
= -2{,}147{,}483{,}520$.
Range: $\texttt{\$s0} \in [-2{,}147{,}483{,}648,\; -2{,}147{,}483{,}521]$.

\noindent\textbf{2.11.3} \texttt{sub \$s0, \$s1, \$s0}: overflow when
$128 - \texttt{\$s0} > 2^{31}-1$ (positive overflow) or
$128 - \texttt{\$s0} < -2^{31}$ (negative overflow).
Positive: $\texttt{\$s0} < 128 - (2^{31}-1) = -2{,}147{,}483{,}519$, so
$\texttt{\$s0} \in [-2{,}147{,}483{,}648,\; -2{,}147{,}483{,}520]$.

%-----------------------------------------------------------------------------
\subsection*{2.12 --- Binary to Instruction Type}
\hypertarget{sol-2.12}{}
\markright{2.12}
\addcontentsline{toc}{subsection}{2.12 --- Binary to Instruction Type}

Binary: \texttt{0000 0010 0001 0100 1000 0000 0010 0000}.

\begin{verbatim}
  000000 10001 01001 00000 00100 000
\end{verbatim}
Wait---let us regroup by field boundaries:
\begin{verbatim}
  000000 00100 10100 10000 00000 100000
  op=0   rs=4  rt=20  rd=16 sh=0  fn=32
\end{verbatim}

Hmm, let me recount the 32 bits:
\begin{verbatim}
  0000 0010 0001 0100 1000 0000 0010 0000
\end{verbatim}
R-type field split (6-5-5-5-5-6):
\begin{verbatim}
  000000 10001 01001 00000 00010 000000
\end{verbatim}
Wait, that does not add up to 32 bits properly either. Let me be careful:
\begin{verbatim}
  Bits: 0000 0010 0001 0100 1000 0000 0010 0000
  Flat: 00000010000101001000000000100000
  R:    000000 10000 10100 10000 00000 100000
        op=0   rs=16 rt=20 rd=16 sh=0  fn=32
\end{verbatim}
op=0, funct=32 $\Rightarrow$ \textbf{R-type, \texttt{add}}.
\texttt{add \$s0, \$s0, \$s4} \quad (rd=16=\$s0, rs=16=\$s0, rt=20=\$s4).

%-----------------------------------------------------------------------------
\subsection*{2.13 --- Instruction to Hex Representation}
\hypertarget{sol-2.13}{}
\markright{2.13}
\addcontentsline{toc}{subsection}{2.13 --- Instruction to Hex Representation}

\texttt{add \$t1, \$s2, \$t1}: R-type.
Registers: rd=\$t1=9, rs=\$s2=18, rt=\$t1=9, shamt=0, funct=32.
\begin{verbatim}
  000000 10010 01001 01001 00000 100000
\end{verbatim}
Hex: \texttt{0x02494820}.

%-----------------------------------------------------------------------------
\subsection*{2.14 --- MIPS Fields to Instruction (R-type)}
\hypertarget{sol-2.14}{}
\markright{2.14}
\addcontentsline{toc}{subsection}{2.14 --- MIPS Fields to Instruction (R-type)}

\texttt{op=0, rs=3, rt=2, rd=3, shamt=0, funct=34}.

op=0, funct=34 $\Rightarrow$ \texttt{sub}.
rs=3=\$v1, rt=2=\$v0, rd=3=\$v1.

Assembly: \texttt{sub \$v1, \$v1, \$v0}.

Binary:
\begin{verbatim}
  000000 00011 00010 00011 00000 100010
\end{verbatim}

%-----------------------------------------------------------------------------
\subsection*{2.15 --- MIPS Fields to Instruction (I-type)}
\hypertarget{sol-2.15}{}
\markright{2.15}
\addcontentsline{toc}{subsection}{2.15 --- MIPS Fields to Instruction (I-type)}

\texttt{op=0x23, rs=1, rt=2, const=0x4}.

op=0x23=35 $\Rightarrow$ \texttt{lw} (load word).
rs=1=\$at, rt=2=\$v0.

Assembly: \texttt{lw \$v0, 4(\$at)}.

Binary:
\begin{verbatim}
  100011 00001 00010 0000000000000100
\end{verbatim}

%-----------------------------------------------------------------------------
\subsection*{2.16 --- Expanding the Register File}
\hypertarget{sol-2.16}{}
\markright{2.16}
\addcontentsline{toc}{subsection}{2.16 --- Expanding the Register File}

128 registers $\Rightarrow$ 7 bits per register field.
4$\times$ instructions $\Rightarrow$ opcode from 6 to 8 bits.

\noindent\textbf{2.16.1} R-type: op(8) + rs(7) + rt(7) + rd(7) + shamt(?)
+ funct(?) = 32 bits. 32 $-$ 8 $-$ 7 $-$ 7 $-$ 7 = 3 bits left for
shamt + funct. The funct field shrinks drastically, limiting the number of
R-type ALU operations.

\noindent\textbf{2.16.2} I-type: op(8) + rs(7) + rt(7) + immediate(?)
= 32 bits. Immediate $= 32 - 8 - 7 - 7 = 10$ bits (reduced from 16).

\noindent\textbf{2.16.3}
\textit{Decrease program size:} more registers means fewer load/store
spills; more opcodes can encode complex operations in one instruction.
\textit{Increase program size:} the smaller immediate field forces more
\texttt{lui}/\texttt{ori} pairs to load constants, and branch offsets become
shorter, possibly requiring extra jumps.

%-----------------------------------------------------------------------------
\subsection*{2.17 --- Bitwise Operations on Registers}
\hypertarget{sol-2.17}{}
\markright{2.17}
\addcontentsline{toc}{subsection}{2.17 --- Bitwise Operations on Registers}

\texttt{\$t0 = 0xAAAAAAAA}, \texttt{\$t1 = 0x12345678}.

\noindent\textbf{2.17.1}
\begin{align*}
  \texttt{sll \$t2, \$t0, 4}: &\quad \texttt{0xAAAAAAAA} \ll 4
    = \texttt{0xAAAAAAA0} \\
  \texttt{or  \$t2, \$t2, \$t1}: &\quad \texttt{0xAAAAAAA0}
    \mid \texttt{0x12345678} = \texttt{0xBABEFEF8}
\end{align*}

\noindent\textbf{2.17.2}
\begin{align*}
  \texttt{srl \$t2, \$t0, 3}: &\quad \texttt{0xAAAAAAAA} \gg 3
    = \texttt{0x15555555} \\
  \texttt{and \$t2, \$t2, \$t1}: &\quad \texttt{0x15555555}
    \mathbin{\&} \texttt{0x12345678} = \texttt{0x10145450}
\end{align*}

\noindent\textbf{2.17.3}
\begin{align*}
  \texttt{srl \$t2, \$t0, 4}: &\quad \texttt{0xAAAAAAAA} \gg 4
    = \texttt{0x0AAAAAAA} \\
  \texttt{andi \$t2, \$t2, -1}: &\quad \text{andi sign-extends the 16-bit
    immediate } \texttt{-1} = \texttt{0xFFFF} \\
  &\quad \texttt{0x0AAAAAAA} \mathbin{\&} \texttt{0x0000FFFF}
    = \texttt{0x0000AAAA}
\end{align*}
Note: \texttt{andi} zero-extends the 16-bit immediate, so \texttt{-1} as a
16-bit value is \texttt{0xFFFF}, which becomes \texttt{0x0000FFFF}.

%-----------------------------------------------------------------------------
\subsection*{2.18 --- Bit Field Extraction}
\hypertarget{sol-2.18}{}
\markright{2.18}
\addcontentsline{toc}{subsection}{2.18 --- Bit Field Extraction}

Extract bits 16:11 from \texttt{\$t0}, place into bits 31:26 of
\texttt{\$t1}:
\begin{verbatim}
    srl  $t2, $t0, 11      # shift bits 16:11 to 5:0
    andi $t2, $t2, 0x3F    # mask to 6 bits
    sll  $t2, $t2, 26      # move to bits 31:26
    lui  $t3, 0x03FF       # $t3 = 0x03FFFFFF
    ori  $t3, $t3, 0xFFFF  # mask for bits 25:0
    and  $t1, $t1, $t3     # clear bits 31:26 of $t1
    or   $t1, $t1, $t2     # insert extracted bits
\end{verbatim}

%-----------------------------------------------------------------------------
\subsection*{2.19 --- Pseudoinstruction Implementation}
\hypertarget{sol-2.19}{}
\markright{2.19}
\addcontentsline{toc}{subsection}{2.19 --- Pseudoinstruction Implementation}

\texttt{not \$t0, \$t1} $\equiv$ bitwise NOT.
\begin{verbatim}
    nor $t0, $t1, $zero    # $t0 = ~($t1 | 0) = ~$t1
\end{verbatim}
One instruction.

%-----------------------------------------------------------------------------
\subsection*{2.20 --- C to MIPS: Array Sum}
\hypertarget{sol-2.20}{}
\markright{2.20}
\addcontentsline{toc}{subsection}{2.20 --- C to MIPS: Array Sum}

\texttt{A = C[0] + C[4];}  with \texttt{A}$=$\texttt{\$t0},
base of \texttt{C}$=$\texttt{\$s0}.
\begin{verbatim}
    lw   $t1, 0($s0)       # $t1 = C[0]
    lw   $t2, 16($s0)      # $t2 = C[4]  (4*4=16)
    add  $t0, $t1, $t2     # A = C[0] + C[4]
\end{verbatim}

%-----------------------------------------------------------------------------
\subsection*{2.21 --- SLT and Branch}
\hypertarget{sol-2.21}{}
\markright{2.21}
\addcontentsline{toc}{subsection}{2.21 --- SLT and Branch}

\texttt{\$s0 = 0x80000000} $= -2^{31}$ (most negative signed value).
\begin{verbatim}
    slt  $t2, $s0, $zero   # -2^31 < 0? Yes -> $t2 = 1
    bne  $t2, $zero, ELSE  # 1 != 0, branch taken
    j    EXIT              # (skipped)
ELSE:
    addi $t2, $t2, 2       # $t2 = 1 + 2 = 3
EXIT:
\end{verbatim}
\textbf{\$t2 = 3.}

%-----------------------------------------------------------------------------
\subsection*{2.22 --- PC Addressing Range}
\hypertarget{sol-2.22}{}
\markright{2.22}
\addcontentsline{toc}{subsection}{2.22 --- PC Addressing Range}

PC $=$ \texttt{0x00400000}.

\noindent\textbf{2.22.1} \texttt{jal} uses a 26-bit word address. The
target address is formed by concatenating the upper 4 bits of PC+4 with the
26-bit field shifted left by 2:
\[
  \text{Range} = [\texttt{0x00000000},\; \texttt{0x0FFFFFFC}]
\]
(any address in the same 256 MB region as PC+4).

\noindent\textbf{2.22.2} \texttt{beq} uses a 16-bit signed offset
(word-aligned):
\[
  \text{Range} = [\text{PC}+4 - 2^{17},\; \text{PC}+4 + 2^{17}-4]
  = [\texttt{0x003E0004},\; \texttt{0x00420000}]
\]
($\pm 128$ KB from PC+4.)

%-----------------------------------------------------------------------------
\subsection*{2.23 --- Proposed \texttt{rpt} Instruction}
\hypertarget{sol-2.23}{}
\markright{2.23}
\addcontentsline{toc}{subsection}{2.23 --- Proposed \texttt{rpt} Instruction}

\textbf{Format:} I-type works well---rs holds the counter register, the
immediate encodes the decrement value, and the offset (or label) points to
the loop target.

\textbf{Issues:} (1) The instruction does two operations (decrement +
conditional branch), violating the MIPS ``simplicity'' design principle.
(2) It writes a register \emph{and} changes the PC in one instruction,
complicating the pipeline. (3) The branch offset may need more bits than
available in a single I-type immediate field.

%-----------------------------------------------------------------------------
\subsection*{2.24 --- MIPS Loop Analysis}
\hypertarget{sol-2.24}{}
\markright{2.24}
\addcontentsline{toc}{subsection}{2.24 --- MIPS Loop Analysis}

\noindent\textbf{2.24.1}
The loop loads \texttt{MemArray[0]} through \texttt{MemArray[9]} and
accumulates the sum in \texttt{\$s2}:
$\texttt{\$s2} = \sum_{k=0}^{9} \texttt{MemArray}[k]$.

\noindent\textbf{2.24.2} C code:
\begin{verbatim}
    int b = 0;
    for (int i = 10; i > 0; i--) {
        b = b + *ptr;   // ptr initially = &MemArray[0]
        ptr++;
    }
\end{verbatim}

\noindent\textbf{2.24.3}
Loop body = 5 instructions, plus 1 initialization = 1 + $5N$
instructions total. The \texttt{bne} in the last iteration falls through.
Total: $5N + 1$.

%-----------------------------------------------------------------------------
\subsection*{2.25 --- C to MIPS: Nested Loop}
\hypertarget{sol-2.25}{}
\markright{2.25}
\addcontentsline{toc}{subsection}{2.25 --- C to MIPS: Nested Loop}

\begin{verbatim}
        add  $t0, $zero, $zero    # i = 0
OUTER:  slt  $t2, $t0, $s0       # i < a?
        beq  $t2, $zero, DONE    # if not, exit
        add  $t1, $zero, $zero   # j = 0
INNER:  slt  $t2, $t1, $s1      # j < b?
        beq  $t2, $zero, ENDIN  # if not, exit inner
        add  $t3, $t0, $t1      # i + j
        sll  $t4, $t1, 4        # 4*j*4 = j*16 (byte offset)
        add  $t4, $s2, $t4      # &D[4*j]
        sw   $t3, 0($t4)        # D[4*j] = i + j
        addi $t1, $t1, 1        # j++
        j    INNER
ENDIN:  addi $t0, $t0, 1        # i++
        j    OUTER
DONE:
\end{verbatim}

%-----------------------------------------------------------------------------
\subsection*{2.26 --- Instruction Count for Nested Loop}
\hypertarget{sol-2.26}{}
\markright{2.26}
\addcontentsline{toc}{subsection}{2.26 --- Instruction Count for Nested Loop}

Static instruction count: 13 instructions (as listed above).

With $a=10$, $b=1$: outer loop runs 10 times, inner loop runs 1 time per
outer iteration.
\begin{itemize}
  \item Outer setup: 1 (init i)
  \item Per outer iteration: 3 (slt, beq, init j) + inner iterations +
        2 (addi i, j OUTER)
  \item Per inner iteration: 7 instructions (slt, beq, add, sll, add, sw,
        addi, j)---actually 8 counting the j INNER.
  \item Final outer exit: 2 (slt, beq)
\end{itemize}
Total $\approx 1 + 10 \times (3 + 1 \times 8 + 2 + 2) + 2
= 1 + 10 \times 15 + 2 = 153$ (approximate; exact count depends on
implementation details).

%-----------------------------------------------------------------------------
\subsection*{2.27 --- MIPS Loop to C}
\hypertarget{sol-2.27}{}
\markright{2.27}
\addcontentsline{toc}{subsection}{2.27 --- MIPS Loop to C}

Reading the assembly:
\begin{verbatim}
    i = 0;
    for (i = 0; i < 10; i++) {
        result = MemArray[i] + i;
    }
\end{verbatim}
Note: \texttt{result} is overwritten each iteration, so the final value
is \texttt{MemArray[9] + 9}.

Actually, looking more carefully at the code: the \texttt{lw} loads from
the current pointer, then adds \texttt{i} to the loaded value and stores
the result in \texttt{\$s2}. The pointer advances by 4 each iteration.

%-----------------------------------------------------------------------------
\subsection*{2.28 --- Optimize MIPS Loop}
\hypertarget{sol-2.28}{}
\markright{2.28}
\addcontentsline{toc}{subsection}{2.28 --- Optimize MIPS Loop}

The variable \texttt{i} is only used for loop control (\texttt{slti \$t0,
\$s1, 10}). We can eliminate it by computing a terminal address instead:
\begin{verbatim}
          addi $t0, $s0, 40       # end = &MemArray[10]
    LOOP: lw   $s2, 0($s0)        # result = *ptr
          addi $s0, $s0, 4        # ptr++
          slt  $t1, $s0, $t0     # ptr < end?
          bne  $t1, $zero, LOOP
\end{verbatim}
Reduced from 6 to 4 instructions per iteration.

%-----------------------------------------------------------------------------
\subsection*{2.29 --- Fibonacci in MIPS}
\hypertarget{sol-2.29}{}
\markright{2.29}
\addcontentsline{toc}{subsection}{2.29 --- Fibonacci in MIPS}

\begin{verbatim}
fib:
    addi $sp, $sp, -16     # allocate stack frame
    sw   $ra, 12($sp)      # save return address
    sw   $a0, 8($sp)       # save n
    sw   $s0, 4($sp)       # save $s0

    slti $t0, $a0, 2       # n < 2?
    beq  $t0, $zero, RECUR # if n >= 2, recurse

    # Base case: return n (n=0 or n=1)
    add  $v0, $a0, $zero
    j    RETURN

RECUR:
    addi $a0, $a0, -1      # arg = n - 1
    jal  fib                # fib(n-1)
    add  $s0, $v0, $zero   # $s0 = fib(n-1)

    lw   $a0, 8($sp)       # restore n
    addi $a0, $a0, -2      # arg = n - 2
    jal  fib                # fib(n-2)

    add  $v0, $s0, $v0     # return fib(n-1)+fib(n-2)

RETURN:
    lw   $s0, 4($sp)       # restore $s0
    lw   $ra, 12($sp)      # restore $ra
    addi $sp, $sp, 16      # free stack frame
    jr   $ra
\end{verbatim}

%-----------------------------------------------------------------------------
\subsection*{2.30 --- Stack Contents for Function Calls}
\hypertarget{sol-2.30}{}
\markright{2.30}
\addcontentsline{toc}{subsection}{2.30 --- Stack Contents for Function Calls}

Each call to \texttt{fib} pushes 16 bytes (4 words): \texttt{\$ra},
\texttt{\$a0} (n), \texttt{\$s0}, and padding. Starting at
\texttt{SP = 0x7FFFFFFC}:

For \texttt{fib(3)}: SP goes to \texttt{0x7FFFFFEC}, then \texttt{fib(2)}
pushes another frame at \texttt{0x7FFFFFDC}, and so on. Each recursive call
adds a 16-byte frame containing the saved \texttt{\$ra}, \texttt{n},
\texttt{\$s0}, and one padding word.

The maximum stack depth for \texttt{fib(n)} is $O(n)$ frames.

%-----------------------------------------------------------------------------
\subsection*{2.31 --- Translate Function to MIPS}
\hypertarget{sol-2.31}{}
\markright{2.31}
\addcontentsline{toc}{subsection}{2.31 --- Translate Function to MIPS}

\begin{verbatim}
f:
    addi $sp, $sp, -16     # allocate frame
    sw   $ra, 12($sp)      # save $ra
    sw   $s0, 8($sp)       # save $s0

    add  $s0, $a2, $a3     # $s0 = c + d (save for later)
    jal  func              # func(a, b) -- $a0,$a1 already set
    add  $a0, $v0, $zero   # first arg = func(a,b)
    add  $a1, $s0, $zero   # second arg = c + d
    jal  func              # func(func(a,b), c+d)

    lw   $s0, 8($sp)       # restore $s0
    lw   $ra, 12($sp)      # restore $ra
    addi $sp, $sp, 16      # free frame
    jr   $ra               # return (result in $v0)
\end{verbatim}

%-----------------------------------------------------------------------------
\subsection*{2.32 --- Tail-Call Optimization}
\hypertarget{sol-2.32}{}
\markright{2.32}
\addcontentsline{toc}{subsection}{2.32 --- Tail-Call Optimization}

\textbf{No}, we cannot use tail-call optimization. The return value of the
first call to \texttt{func(a,b)} must be saved before making the second call
to \texttt{func}. The function \texttt{f} performs work \emph{between} the
two calls, so the last action is not a single function call---it depends on
the result of a prior call.

%-----------------------------------------------------------------------------
\subsection*{2.33 --- Stack State Before Return}
\hypertarget{sol-2.33}{}
\markright{2.33}
\addcontentsline{toc}{subsection}{2.33 --- Stack State Before Return}

Before \texttt{f} returns: \texttt{\$s0} is restored to its original
(caller-saved) value. \texttt{\$s1}, \texttt{\$s3}, \texttt{\$s4} were never
modified by \texttt{f}, so they hold whatever value the caller set.
\texttt{\$t5} through \texttt{\$t7} are caller-saved---\texttt{func} may
have changed them, so their values are \textbf{unknown}.

%-----------------------------------------------------------------------------
\subsection*{2.34 --- ASCII String to Integer}
\hypertarget{sol-2.34}{}
\markright{2.34}
\addcontentsline{toc}{subsection}{2.34 --- ASCII String to Integer}

\begin{verbatim}
atoi:
    add  $v0, $zero, $zero  # result = 0
    add  $t1, $zero, $zero  # sign = 0 (positive)
    lb   $t0, 0($a0)        # load first char
    li   $t3, 0x2D           # '-'
    bne  $t0, $t3, LOOP     # not minus, skip
    li   $t1, 1              # sign = 1 (negative)
    addi $a0, $a0, 1         # advance past '-'
LOOP:
    lb   $t0, 0($a0)        # load char
    beq  $t0, $zero, DONE   # null terminator -> done
    slti $t2, $t0, 48       # < '0'?
    bne  $t2, $zero, ERR
    slti $t2, $t0, 58       # < '9'+1?
    beq  $t2, $zero, ERR
    addi $t0, $t0, -48      # digit = char - '0'
    li   $t2, 10
    mult $v0, $t2
    mflo $v0                 # result *= 10
    add  $v0, $v0, $t0      # result += digit
    addi $a0, $a0, 1
    j    LOOP
DONE:
    beq  $t1, $zero, POS
    sub  $v0, $zero, $v0    # negate if negative
POS:
    jr   $ra
ERR:
    li   $v0, -1
    jr   $ra
\end{verbatim}

%-----------------------------------------------------------------------------
\subsection*{2.35 --- Memory Store and Endianness}
\hypertarget{sol-2.35}{}
\markright{2.35}
\addcontentsline{toc}{subsection}{2.35 --- Memory Store and Endianness}

\texttt{lbu \$t0, 0(\$t1)}: loads byte at \texttt{0x10000000} into
\texttt{\$t0} (zero-extended to 32 bits).

\texttt{sw \$t0, 0(\$t1)}: stores the 32-bit value in \texttt{\$t0} back
to \texttt{0x10000000}.

\noindent\textbf{2.35.1} Big-endian: the byte was at address 0 of the word
(MSB). After \texttt{lbu}, \texttt{\$t0 = 0x000000XX} where \texttt{XX} is
the byte. After \texttt{sw}, the word at \texttt{0x10000000} becomes
\texttt{0x000000XX}---only the original MSB remains.

\noindent\textbf{2.35.2} Little-endian: the byte at address 0 is the LSB.
After \texttt{lbu}, \texttt{\$t0 = 0x000000XX}. After \texttt{sw}, the word
becomes \texttt{0x000000XX}---same result, but the LSB position is preserved.

%-----------------------------------------------------------------------------
\subsection*{2.36 --- Create 32-bit Constant}
\hypertarget{sol-2.36}{}
\markright{2.36}
\addcontentsline{toc}{subsection}{2.36 --- Create 32-bit Constant}

\texttt{0x0000FACE}: upper 16 bits are 0, lower 16 are \texttt{0xFACE}.
\begin{verbatim}
    ori $t1, $zero, 0xFACE   # $t1 = 0x0000FACE
\end{verbatim}
One instruction suffices because the upper half is zero and \texttt{ori}
zero-extends its 16-bit immediate.

%-----------------------------------------------------------------------------
\subsection*{2.37 --- Atomic Set-Max with \texttt{ll/sc}}
\hypertarget{sol-2.37}{}
\markright{2.37}
\addcontentsline{toc}{subsection}{2.37 --- Atomic Set-Max with \texttt{ll/sc}}

\begin{verbatim}
setmax:
    # $a0 = x (address), $a1 = s (value)
RETRY:
    ll   $t0, 0($a0)       # $t0 = *x (linked)
    slt  $t1, $t0, $a1     # *x < s?
    beq  $t1, $zero, SKIP  # if *x >= s, skip
    sc   $a1, 0($a0)       # try *x = s
    beq  $a1, $zero, RETRY # if sc failed, retry
SKIP:
    jr   $ra
\end{verbatim}

Note: after \texttt{sc}, \texttt{\$a1} is overwritten with success/failure.
If the store fails (another processor modified the location), we retry.

%-----------------------------------------------------------------------------
\subsection*{2.38 --- Concurrent Critical Section}
\hypertarget{sol-2.38}{}
\markright{2.38}
\addcontentsline{toc}{subsection}{2.38 --- Concurrent Critical Section}

When two processors execute simultaneously:
\begin{enumerate}
  \item Both execute \texttt{ll} and read the same value of \texttt{*x}.
  \item Both compare and decide to update.
  \item One processor's \texttt{sc} succeeds (returns 1); the other's fails
        (returns 0) because the \texttt{ll/sc} pair detects the intervening
        write.
  \item The failing processor retries from the \texttt{ll}, now reading the
        updated value.
\end{enumerate}
This guarantees atomicity---only one processor's update takes effect per
\texttt{ll/sc} pair.

%-----------------------------------------------------------------------------
\subsection*{2.39 --- New Arithmetic Instructions}
\hypertarget{sol-2.39}{}
\markright{2.39}
\addcontentsline{toc}{subsection}{2.39 --- New Arithmetic Instructions}

CPI: arith=1, L/S=10, branch=3.
IC: arith=500M, L/S=300M, branch=100M.

Original CPU time (relative):
$500 \times 1 + 300 \times 10 + 100 \times 3 = 500 + 3000 + 300 = 3800$
(in millions of cycles).

\noindent\textbf{2.39.1}
New arith IC $= 0.75 \times 500 = 375$M. Clock cycle time $\times 1.1$.
\[
  \text{New cycles} = 375 \times 1 + 300 \times 10 + 100 \times 3
  = 375 + 3000 + 300 = 3675
\]
\[
  \text{New time} = 3675 \times 1.1 = 4042.5
  \quad\text{vs.}\quad 3800
\]
\textbf{No}, this is a bad design choice---execution time increases by 6.4\%.

\phantomsection\addcontentsline{toc}{subsubsection}{2.39.2}
\medskip\noindent\textbf{2.39.2}
Double arithmetic performance $\Rightarrow$ arith CPI $= 0.5$:
$\text{New cycles} = 500(0.5) + 3000 + 300 = 3550$.
Speedup $= 3800/3550 = 1.070$ ($7.0\%$).

10$\times$ arithmetic: arith CPI $= 0.1$:
$\text{New cycles} = 500(0.1) + 3000 + 300 = 3350$.
Speedup $= 3800/3350 = 1.134$ ($13.4\%$).

Even 10$\times$ improvement in arithmetic yields only 13\% overall speedup
(Amdahl's law---arithmetic is only $500/3800 = 13\%$ of total time).

%-----------------------------------------------------------------------------
\subsection*{2.40 --- CPI and Instruction Mix}
\hypertarget{sol-2.40}{}
\markright{2.40}
\addcontentsline{toc}{subsection}{2.40 --- CPI and Instruction Mix}

70\% arithmetic (2 cyc), 10\% L/S (6 cyc), 20\% branch (3 cyc).

\noindent\textbf{2.40.1}
$\text{CPI} = 0.7(2) + 0.1(6) + 0.2(3) = 1.4 + 0.6 + 0.6 = 2.6$.

\noindent\textbf{2.40.2}
25\% improvement $\Rightarrow$ new CPI $= 2.6/1.25 = 2.08$.
$0.7 \cdot x + 0.6 + 0.6 = 2.08 \Rightarrow 0.7x = 0.88
\Rightarrow x = 1.257$ cycles per arithmetic instruction.

\noindent\textbf{2.40.3}
50\% improvement $\Rightarrow$ new CPI $= 2.6/1.5 = 1.733$.
$0.7 \cdot x + 0.6 + 0.6 = 1.733 \Rightarrow 0.7x = 0.533
\Rightarrow x = 0.762$ cycles.

This requires arithmetic CPI $< 1$, which is physically unrealistic for
a single-issue processor. \textbf{Cannot be achieved} by improving only
arithmetic instructions.

%-----------------------------------------------------------------------------
\subsection*{2.41 --- Scaled Offset for Exercise 2.4}
\hypertarget{sol-2.41}{}
\markright{2.41}
\addcontentsline{toc}{subsection}{2.41 --- Scaled Offset for Exercise 2.4}

With scaled offset addressing, \texttt{lw \$s0, A(\$s0*4)} would combine
the shift-left-by-2 and the base addition into the addressing mode. This
eliminates the explicit \texttt{sll} and \texttt{add} instructions. The
9-instruction sequence in 2.4 could be reduced to roughly 5--6 instructions.

%-----------------------------------------------------------------------------
\subsection*{2.42 --- Scaled Offset for Exercise 2.7}
\hypertarget{sol-2.42}{}
\markright{2.42}
\addcontentsline{toc}{subsection}{2.42 --- Scaled Offset for Exercise 2.7}

Similarly, with scaled offset loads, the \texttt{sll} (shift by 3 for
8-byte elements) and \texttt{add} to form the address can be folded into a
single load instruction like \texttt{lw \$t0, A(\$s3*8)}. This reduces the
8-instruction sequence from 2.7 by 2--3 instructions.
